% ID: MAT-GAN-RET-002
% Titolo: Modello lineare e scarti su dati simulati
% Disciplina: Matematica
% Area: geometria_analitica
% Argomento: retta
% Classe: Prima o terza secondaria di secondo grado
% Difficolta: 3
% Tipo: modellizzazione
% Tag: matematica, geometria_analitica, retta
% Fonte: laboratorio_modello_lineare_anonimizzato
% Provenienza: materiale fornito dal docente; testo anonimizzato; soluzione redatta per Exergo

\begin{esercizio}
Ai tempi $t=0,2,4,6,8,10,12$ secondi corrispondono le posizioni simulate
$s=0{,}2;1{,}4;2{,}5;3{,}8;5{,}0;6{,}1;7{,}4$ metri.
Costruisci la retta che unisce il primo e l'ultimo punto e calcola gli scarti dei dati dalla retta.
Confronta graficamente dati e modello; prevedi la posizione a 15 secondi.
Sostituisci temporaneamente 3,8 con 9,8: il coefficiente angolare della retta cambia?
\end{esercizio}

\begin{soluzione}
Il coefficiente angolare è $(7{,}4-0{,}2)/12=0{,}6$, quindi $s_m(t)=0{,}2+0{,}6t$.
Le previsioni ai tempi assegnati sono $0{,}2;1{,}4;2{,}6;3{,}8;5{,}0;6{,}2;7{,}4$;
gli scarti dato meno modello sono $0;0;-0{,}1;0;0;-0{,}1;0$.
A 15 secondi il modello dà $9{,}2$ metri: è un'estrapolazione, non una misura.
Il punto alterato ha scarto 6 metri ma non cambia i punti estremi, quindi la pendenza rimane $0{,}6$.
Nel grafico XY i dati simulati si confrontano con la retta; uno scarto negativo
indica un dato sotto la retta, non automaticamente una misura sbagliata.
Le pendenze successive originali sono $0{,}60;0{,}55;0{,}65;0{,}60;0{,}55;0{,}65$;
con il dato alterato le due pendenze adiacenti diventano $3{,}65$ e $-2{,}40$.
\end{soluzione}
